% INSTRUCTOR'S HANDWRITTEN DERIVATION attached to the AB-effect/monopole pages.
% Placement: immediately following the Aharonov--Bohm effect, before monopoles.
\section{Berry Phase}

The Aharonov--Bohm effect showed that the relative phase of a quantum amplitude can depend on the geometry of a path even where the magnetic field vanishes. A related geometric phase appears when a Hamiltonian depends on external parameters that vary slowly along a path. In this case the path lies in \emph{parameter space}, rather than necessarily in the particle's physical position space.

\subsection*{1) Instantaneous eigenstates and adiabatic evolution}

Let a Hamiltonian depend smoothly on real control parameters $\vect R=(R_1,\ldots,R_d)$:
\[\op H(\vect R)\ket{n(\vect R)}=E_n(\vect R)\ket{n(\vect R)},\qquad
\braket{n(\vect R)}{n(\vect R)}=1.\]
Assume that $\vect R(t)$ changes sufficiently slowly and that the chosen eigenvalue remains nondegenerate and separated by a gap from the other energy levels along the path. Under the adiabatic approximation a state prepared in $\ket{n(\vect R(0))}$ follows the corresponding instantaneous eigenspace. Its evolution contains a dynamical phase and a second, geometrical phase:
\[\boxed{\ket{\psi(t)}\simeq
 e^{-\frac\ii\hbar\int_0^t E_n(\vect R(t'))\,\dd t'}
 e^{\ii\gamma_n(t)}\ket{n(\vect R(t))}.}\]
We determine $\gamma_n$ from the Schr\"odinger equation, rather than assuming its form.

Write $\ket\psi=e^{-\ii\int E_n\dd t/\hbar}e^{\ii\gamma_n}\ket n$ and substitute into
\[\ii\hbar\frac{\dd}{\dd t}\ket\psi=\op H(\vect R(t))\ket\psi.\]
Differentiation gives the identity within the adiabatic approximation,
\[\left(E_n-\hbar\dot\gamma_n\right)\ket n
   +\ii\hbar\frac{\dd}{\dd t}\ket n=E_n\ket n.\]
Project onto the instantaneous normalized eigenstate to find
\[\boxed{\dot\gamma_n(t)=\ii\bra{n(\vect R(t))}
\frac{\dd}{\dd t}\ket{n(\vect R(t))}.}\]
The normalization condition implies that $\braket n{\dot n}$ is purely imaginary, so $\dot\gamma_n$ is real. Apply the chain rule,
\[\frac{\dd}{\dd t}\ket{n(\vect R(t))}
 =\sum_\mu\frac{\partial}{\partial R_\mu}\ket{n(\vect R)}\dot R_\mu,\]
to rewrite the accumulated phase as an integral along the \emph{parameter path} $C$:
\[\boxed{\gamma_n(C)=\int_C\vect{\mathcal A}_n(\vect R)\cdot\dd\vect R,
\qquad\vect{\mathcal A}_n(\vect R)
 =\ii\braket{n(\vect R)}{\vect\nabla_{\!R}n(\vect R)}.}\]
The vector field $\vect{\mathcal A}_n$ is the \emph{Berry connection}. It depends on how the phase of the instantaneous eigenstate has been chosen.

\subsection*{2) Gauge transformations of the eigenstate}

The normalized eigenket is physically unchanged by a parameter-dependent phase. Choose a new convention
\[\ket{n'(\vect R)}=e^{-\ii\beta(\vect R)}\ket{n(\vect R)},\qquad
 \beta(\vect R)\in\mathbb R.\]
Substituting it into the Berry connection shows explicitly how the additional derivative acts:
\[\begin{aligned}
\vect{\mathcal A}'_n
 &=\ii\bra n e^{\ii\beta}\vect\nabla_{\!R}
    \bigl(e^{-\ii\beta}\ket n\bigr)\\
 &=\ii\braket n{\vect\nabla_{\!R}n}
       +\vect\nabla_{\!R}\beta\\
 &=\boxed{\vect{\mathcal A}_n+\vect\nabla_{\!R}\beta}.
\end{aligned}\]
Hence an \emph{open-path} integral transforms as
\[\gamma_n'(C)=\gamma_n(C)+\beta(\vect R_f)-\beta(\vect R_i).\]
For a closed path $C$, with a single-valued choice of phase along the loop, the two endpoint contributions cancel modulo $2\pi$. Thus the Berry phase around a closed cycle is a gauge-invariant quantity:
\[\boxed{e^{\ii\gamma_n(C)}\quad\text{is gauge invariant for a closed }C.}\]
The distinction matters: neither the connection at one point nor the Berry phase assigned to an open segment is independently gauge invariant. A physically measured phase comes from interference, from a cyclic state evolution, or from an appropriate endpoint comparison.

\subsection*{3) Berry curvature and geometry of parameter space}

For three control parameters define the Berry curvature by
\[\boxed{\vect{\mathcal F}_n(\vect R)
 =\vect\nabla_{\!R}\times\vect{\mathcal A}_n(\vect R).}\]
It is unaffected by the gauge change because $\vect\nabla_{\!R}\times\vect\nabla_{\!R}\beta=0$. For a surface $\Sigma$ whose boundary is $C$, Stokes' theorem gives
\[\boxed{\gamma_n(C)=\oint_C\vect{\mathcal A}_n\cdot\dd\vect R
=\int_\Sigma\vect{\mathcal F}_n\cdot\dd\vect S\quad(\mathrm{mod}\ 2\pi),}\]
provided the eigenstate can be chosen smoothly over the surface; singular gauges or degeneracies require appropriate patches. In more than three parameters use the antisymmetric tensor
\[\mathcal F_{n,\mu\nu}
 =\partial_\mu\mathcal A_{n,\nu}-\partial_\nu\mathcal A_{n,\mu}
 =\ii\left(\braket{\partial_\mu n}{\partial_\nu n}
           -\braket{\partial_\nu n}{\partial_\mu n}\right).\]

\begin{center}
\begin{tikzpicture}[>=Latex,scale=0.85]
 \draw[->] (-0.3,0)--(5.6,0) node[right] {$R_1$};
 \draw[->] (0,-0.3)--(0,3.0) node[above] {$R_2$};
 \fill[blue!8] (2.7,1.3) ellipse (1.6 and .96);
 \draw[very thick,blue!70!black,postaction={decorate},
   decoration={markings,mark=at position .12 with {\arrow{>}}}]
 (2.7,1.3) ellipse (1.6 and .96);
 \node[blue!70!black] at (4.4,2.3) {$C$};
 \node at (2.65,1.2) {$\Sigma$};
 \draw[->,thick,red!75!black] (3,1.6)--(3.38,2.0);
 \node[red!75!black,anchor=west] at (3.42,2.04) {$\vect{\mathcal F}_n$};
\end{tikzpicture}
\end{center}

The analogy with the AB effect is now precise: the AB phase involves the electromagnetic vector potential integrated over a path in physical space; the Berry phase involves a connection generated by the \emph{parameter dependence of eigenstates}. Their common mathematical structure is a line integral of a connection, its gauge transformation by a gradient, and a gauge-invariant curvature or closed-loop phase. The Berry curvature is a property of the quantum eigenstate bundle, not an electromagnetic field in ordinary space. This construction will reappear when Bloch-band eigenstates depend on crystal momentum.